\chapter[Analysis Tools]{How to Analysis}
This chapter starts with an overview on imaging and how that is performed. Then we talk a bit about how we actually want to measure the flux of the minihalos.
Could combine with actual results.

\section{Final datasets}
\begin{itemize}
    \item Could have a table here talking about each minihalo, how much data is remaining, spws that are flagged out, some sort of visual. 
    \item Maybe bar graph of spws, line or point graph for minihalo C/D (either total or differentiated by style) normalised to [0, 1]
\end{itemize}

\section{Calibration Error}
\begin{itemize}
    \item Calibration error exists, and sources can vary in their flux amount.
    \item We can try and constrain the error and variance by comparing the secondary calibrator models for D config and the point source fluxes for the images.
    \item Fluxes for secondaries can vary a lot over a year (100\%?) maybe, but lower per month. Can then compare to differences in sources in the image.
    \item If a trend is observed, likely calibration error. If no trend or close models, constrain calibration error and can see the variance in AGN etc.
    \item Can retrieve the point source flux by cutting out diffuse sources (uvrange) and by masking to a tight border.
\end{itemize}

\section{Calculating the minihalo flux}
\begin{itemize}
    \item We can do this in a variety of ways, but we don't want the AGN contamination so we should subtract this out.
    \item Do this per dataset to account for variation over time, incorrect subtractions can result in negative bowls in emission, or remaining flux.
    \item Other problematic sources are dealt with the same way.
    \item Image the subtracted data and then we create a contour for significant emission (3 sigma) from the residual RMS.
    \item Calculate total flux in the region.
    \item Then we calculate the error using this formula from G14 modified for the AGN error since we think its lower.
\end{itemize}

